\section[Repository]{Preparing the repository}
\sectionframe{03}{Preparing the repository}
  {What can the repository tell the agent, so that we do not have to?}
  {\faFolderOpen}

% ------------------------------------------------------ the layout and its files
% Left: a standard layout, adapted from Cookiecutter Data Science, with one AGENTS.md. Right: the
% same repository later, its AGENTS.md split into three files. One picture, so that the arrow can
% run from one tree to the other: the rows are the theme's folder-tree rows, the pages are drawn
% here. Both trees start at the same height with a row pitch of 3.5 mm; the arrow leaves the
% AGENTS.md row on the left (row 9, y = -12.6 - 8 * 3.5) for rows 5 to 7 on the right.
\begin{frame}[source={DrivenData (2026), Cookiecutter Data Science}]{Setting up the repository}
  \begin{tikzpicture}[x=1mm,y=1mm,
      head/.style={anchor=north west,inner sep=0pt,font=\small\bfseries,text=udsdark},
      page/.style={draw=udsgreylight,fill=white,line width=0.5pt,rounded corners=1.2mm},
      split/.style={udsblue,line width=0.9pt},
      splitarrow/.style={split,-{Stealth[length=2mm,width=1.6mm]}}]
    \useasboundingbox (0,0) rectangle (144,-55);
    % ---- left: the layout, and a link to the template it follows
    \node[head] at (0,0) {A standard layout};
    \node[anchor=north west,inner sep=0pt,font=\footnotesize] at (0,-4.6)
      {\href{https://cookiecutter-data-science.drivendata.org/\#directory-structure}%
        {\faGlobe\enspace\underline{\fn{cookiecutter-data-science}}}};
    \draw[page] (0,-9.4) rectangle (60,-54.4);
    \begin{scope}[shift={(2.4,-12.6)}]
      \treestart{note at=30,row height=3.5}
      \tdir{0}{sem-seg/}{}
      \tdir{1}{data/raw/}{never edited}
      \tdir{1}{data/processed/}{made by the code}
      \tdir{1}{models/}{trained models}
      \tdir{1}{notebooks/}{exploring}
      \tdir{1}{reports/figures/}{figures for papers}
      \tdir{1}{src/sem\_seg/}{the code}
      \tdir{1}{tests/}{}
      \tfile{1}{\textcolor{udsblue}{\textbf{AGENTS.md}}}{\textcolor{udsblue}{for the agent}}
      \tfile{1}{environment.yml}{}
      \tfile{1}{pyproject.toml}{}
      \tfile{1}{README.md}{for people}
    \end{scope}
    % ---- right: the same repository once AGENTS.md has grown; what stays the same in grey
    \node[head] at (74,0) {As the project grows};
    \draw[page] (74,-9.4) rectangle (144,-40.3);
    \begin{scope}[shift={(76.4,-12.6)}]
      \treestart{note at=33.6,row height=3.5}
      \tdir{0}{sem-seg/}{}
      \tdirgrey{1}{data/ \dots}{as before}
      \tdirgrey{1}{models/ \dots\ reports/}{}
      \tdirgrey{1}{src/sem\_seg/, tests/}{}
      \tfile{1}{\textcolor{udsblue}{\textbf{AGENTS.md}}}{commands, rules, pointers}
      \tfile{1}{GLOSSARY.md}{the project's words}
      \tfile{1}{CODING\_STANDARDS.md}{how code is written here}
      \tfilegrey{1}{environment.yml, \dots}{}
    \end{scope}
    \node[anchor=north west,inner sep=0pt,font=\footnotesize,text=udsdark,text width=70mm,
      align=left] at (74,-43.4) {\fn{AGENTS.md} keeps what every session needs and points to
      the two new files};
    % ---- one file becomes three: from the AGENTS.md row on the left to rows 5 to 7 on the right
    \draw[split] (49.5,-40.6) -- (62,-40.6);% starts clear of the note "for the agent"
    \foreach \y in {-26.6,-30.1,-33.6} \draw[splitarrow] (62,-40.6) to[out=0,in=180] (74,\y);
  \end{tikzpicture}
  \note{\textbf{9:50--10:50}\par
    Before any file for the agent: a layout that people and agents already know.\par
    This one follows Cookiecutter Data Science, a template for data projects, with the code in
      src and its tests beside it: raw data that is never edited, processed data that the code
      makes, models, notebooks for exploring, figures for papers.\par
    An agent finds its way in it the way a new colleague does.\par
    At the root sits AGENTS.md, the one file written for the agent.\par
    As the project grows, AGENTS.md keeps what every session needs, and two files take the rest:
      GLOSSARY.md for the project's words, CODING\_STANDARDS.md for how code is written here.}
\end{frame}

% --------------------------------------------------------------- AGENTS.md
% The example file beside what such a file should hold and leave out, after the prompting guides
% of Anthropic and OpenAI (SOURCES.md)
\begin{frame}[fragile,source={Anthropic (2026), Prompting Claude Opus 5.5 · OpenAI (2026), Using GPT-6}]{\texttt{AGENTS.md}}
  \begin{columns}[T,onlytextwidth]
    \begin{column}{0.57\textwidth}
\begin{codefile}[basicstyle=\ttfamily\fontsize{7.2}{8.75}\selectfont]{AGENTS.md}
# sem-seg
## Setup and checks
- Environment: conda env `sem-seg`, activated beforehand
- Fast check (CPU, synthetic data): `pytest -q`
- Lint, types: `ruff check`, `mypy src` §\textcolor{udsred}{\sffamily\scriptsize\faArrowLeft\ after every edit}§
## Conventions
- Terms: `GLOSSARY.md` (specimen, micrograph, patch) §\smash{\textcolor{udsblue}{\faLink}}§
- Decisions: `docs/adr/` §\smash{\textcolor{udsblue}{\faLink}}§
- Lengths in micrometres; pixel size from metadata
## Boundaries
- Never write to `data/raw/`; use `data/processed/`
- No cluster jobs (`sbatch`): the queue is shared
- Ask before adding a dependency
## When you finish
- Report changed files, checks run, open assumptions
\end{codefile}
    \end{column}
    \begin{column}{0.40\textwidth}
      \begin{tikzpicture}[x=1mm,y=1mm,
          head/.style={anchor=north west,inner sep=0pt,font=\footnotesize\bfseries},
          item/.style={anchor=north west,inner sep=0pt,font=\footnotesize,text=udsdark},
          why/.style={anchor=north west,inner sep=0pt,font=\scriptsize,text=udsgrey}]
        \useasboundingbox (0,0) rectangle (57,-54);
        \node[head,text=testgreen] at (0,0) {Put in};
        \foreach \t [count=\i] in {{commands to set up, test and check},
            pointers to what the code cannot show,what needs your approval,
            {when it is done, and what to report},the reason behind a rule}
          \node[item] at (0,-0.7-\i*4.3) {\textcolor{testgreen}{\faCheck}\enspace\t};
        \node[head,text=udsred] at (0,-27.8) {Leave out};
        \foreach \t/\w [count=\i] in {{rules that clash or are out of date}/{the agent may stop and wait},
            {vague wishes: ``write clean code''}/{name the exact mistake instead},
            {``think carefully'', ``be thorough''}/{current models already do}} {
          \node[item] at (0,-25-\i*7.6) {\textcolor{udsred}{\faTimes}\enspace\t};
          \node[why] at (4.3,-28.8-\i*7.6) {\w};
        }
      \end{tikzpicture}
    \end{column}
  \end{columns}
  \note{\textbf{10:50--12:10}\par
    AGENTS.md is the one file that tools supporting it load at the start of every session.\par
    It holds what the agent cannot find out alone: the commands to set up, test and check; what
      the code cannot show, or which file shows it; what needs our approval; when the work is done
      and what to report.\par
    Fifteen lines, each preventing a specific mistake, and one says why: the cluster queue is
      shared.\par
    The checks of section 2 are here, so the agent runs them after every edit; the boundaries are
      requests, not locks.\par
    Today's models follow such files closely: OpenAI notes that its latest model is more
      sensitive to files such as AGENTS.md, and that unclear or conflicting guidance can make it
      stop early.\par
    So leave out rules that clash or are out of date; vague wishes such as write clean code,
      where naming the exact mistake works better; and lines such as think carefully, which
      Anthropic advises removing from system prompts, since the model decides how much to
      think.\par
    Keep it short, and review it like code.}
\end{frame}

% ---------------------------------------------------- glossary, standards
% \tagchip{text}: a tag naming what a file holds
\newcommand\tagchip[1]{\tikz[baseline=(t.base)]\node[draw=udsblue,fill=bluetint,line width=0.4pt,
  rounded corners=0.6mm,inner xsep=1mm,inner ysep=0.5mm,font=\scriptsize,text=udsdark] (t) {#1};}
\begin{frame}[fragile,source={Pocock (2026), mattpocock/skills v1.3.1: domain-modeling, code-review, retro}]{\texttt{GLOSSARY.md} and \texttt{CODING\_STANDARDS.md}}
  \begin{columns}[T,onlytextwidth]
    \begin{column}{0.485\textwidth}
      {\scriptsize\color{udsgrey}holds\enspace}\tagchip{term and meaning}\enspace
        \tagchip{words to avoid}\enspace\tagchip{code name}\par\vskip1.6mm
\begin{codefile}[basicstyle=\ttfamily\fontsize{7}{8.5}\selectfont]{GLOSSARY.md\enspace\textcolor{udsgreymid}{\sffamily\scriptsize read when pointed to}}
# Glossary
Specimen: one physical sample (`specimen_id`).
  Avoid: §\textcolor{udsred}{\bfseries sample}§, coupon. §\textcolor{udsred}{\sffamily\scriptsize\faArrowLeft\ the word in our request}§
Area: one imaged region of a specimen.
Micrograph: one SEM image of an area
  (`micrograph_id`). Avoid: image, frame.
Patch: a crop of a micrograph used in
  training (`patch_id`). Avoid: tile.
Fold: specimens held out together; every
  specimen belongs to exactly one fold.
\end{codefile}
    \end{column}
    \begin{column}{0.485\textwidth}
      {\scriptsize\color{udsgrey}holds\enspace}\tagchip{seeds}\enspace\tagchip{input checks}\enspace
        \tagchip{tests}\enspace\tagchip{paths}\par\vskip1.6mm
\begin{codefile}[basicstyle=\ttfamily\fontsize{7}{8.5}\selectfont]{CODING\_STANDARDS.md\enspace\textcolor{udsgreymid}{\sffamily\scriptsize read in review}}
# Coding standards
- Every random step takes a `seed`; the
  same seed gives the same result.
- Public functions check their inputs and
  raise `ValueError` naming the problem.
- Tests call public functions on small,
  hand-written fixtures, never real data.
- Expected values come from the fixture,
  never recomputed the way the code does.
- Paths come from `config.py`, never
  written into a function.
\end{codefile}
    \end{column}
  \end{columns}
  \note{\textbf{12:10--13:30}\par
    Two files the agent reads only when something points to them.\par
    A glossary holds the project's terms: what each one means, the words to avoid and the name the
      code uses for it; our request said sample, and here sample is a word to avoid.\par
    The clarifying step of section 7 writes terms into it as soon as they are agreed.\par
    Coding standards hold the rules a reviewer checks: seeds, input checks, how tests are built,
      where paths come from; a fixture is the small, hand-made input a test runs on.\par
    In the workflow of section 7 the review step reads it, not the implementing agent, which
      already carries the most context.\par
    Rules that a linter can check belong in Ruff's settings instead, so that every run checks
      them.}
\end{frame}
